% Image credits for the photographs and AI illustrations of Book 5 (University, Year 3).
% Machine-readable license records live in images/book5/CREDITS.md; keep the
% two files in sync. The Book 5 writer owns this file: add one \item per
% photograph (inside the itemize, by chapter) as photographs land.
\chapter*{图片致谢}

本书中的插图、地图和图表均为原创；基于已发表数据绘制的地图和图表，在图注或本列表中注明来源。照片依各自的自由许可使用，谨向作者致谢；每张照片的完整来源链接和许可条款网址，均列于本书代码库中的\texttt{images/book5/CREDITS.md}。

\bigskip
除照片和原创插图外，部分图为 AI 生成的插图，使用 OpenAI 图像生成功能，依据本书编辑撰写的提示词制作，并在纳入前经过科学准确性审阅。这些图只呈现通用场景，从不描绘真实地点、真实人物或真实标本。每幅生成图像的完整提示词记录于代码库中的\texttt{images/book5/ai/PROMPTS.md}。

\bigskip
\noindent 照片与数据，按章列出：
\begin{itemize}\small
  \item 第~1~章（地震学与地球成像）：PREM，A.~M.~Dziewonski 与 D.~L.~Anderson（1981），取自 ObsPy 中的表格；ak135，B.~L.~N.~Kennett、E.~R.~Engdahl 与 R.~Buland（1995）。
  \item 第~3~章（重力与地球形状）：孪生的 GRACE Follow-On 卫星，概念图，NASA/JPL-Caltech，公有领域（维基共享资源）。冰盖质量平衡数据来自 IMBIE 团队（2021），英国极地数据中心，开放政府许可证 v3.0。
  \item 第~4~章（地磁学与古地磁学）：国际地磁参考场第十四代，IAGA 第 V 分部工作组 V-MOD，系数由 NOAA NCEI 发布。
  \item 第~5~章（岩石圈热模型）：海底年龄，M.~Seton 等人（2020），EarthByte，CC BY 4.0；ETOPO1 全球地形数据，NOAA NCEI，公有领域。
  \item 第~6~章（地幔对流与板块驱动力）：Ballyhillin 的抬升海滩，Malin Head，A.~F.~Borchert，CC BY-SA 4.0（维基共享资源）；PREM，A.~M.~Dziewonski 与 D.~L.~Anderson（1981）。
  \item 第~7~章（变形机制与剪切带）：眼状糜棱岩，Mikenorton，CC BY-SA 4.0（维基共享资源）。
  \item 第~9~章（伸展构造与走滑构造）：死亡谷以南的盆岭地形，Ken Lund，CC BY-SA 2.0（维基共享资源）。
  \item 第~14~章（稳定同位素地球化学）：LR04 底栖有孔虫 $\delta^{18}$O 叠加记录，L.~E.~Lisiecki 与 M.~E.~Raymo（2005）。
  \item 第~15~章（水–岩相互作用）：Mammoth 温泉国家公园 Palette Spring 的钙华阶地，Jim Peaco，国家公园管理局，公有领域（维基共享资源）。平衡常数取自 PHREEQC 数据库 phreeqc.dat，D.~L.~Parkhurst 与 C.~A.~J.~Appelo，美国地质调查局，公有领域。
  \item 第~16~章（层序地层学）：犹他州 Helper 以上的 Book Cliffs，David Mackie，公有领域（维基共享资源）；LR04 底栖有孔虫 $\delta^{18}$O 叠加记录，L.~E.~Lisiecki 与 M.~E.~Raymo（2005）。
  \item 第~17~章（相模型）：Arisaig 组中的丘状交错层理，Michael C.~Rygel，CC BY-SA 3.0（维基共享资源），已裁切。
  \item 第~18~章（定量地貌学）：尼亚加拉峡谷与马蹄瀑布，Ken Lund，CC BY-SA 2.0（维基共享资源），已裁切。
  \item 第~19~章（地球流体动力学）：NCEP--NCAR 再分析第 1 版的纬向平均风场，数据由美国科罗拉多州博尔德市 NOAA PSL 提供；世界海洋图集 2023，NOAA NCEI。
  \item 第~20~章（大气动力学）：冰岛附近的一个低压系统，NASA/GSFC，MODIS 快速响应团队，Jacques Descloitres，公有领域（维基共享资源）；NCEP--NCAR 再分析第 1 版的纬向平均风场，数据由美国科罗拉多州博尔德市 NOAA PSL 提供。
  \item 第~21~章（海洋动力学）：从科考船投放的 Argo 浮标，NOAA 全球海洋监测与观测计划，公有领域（维基共享资源）。
  \item 第~22~章（气候模式）：美国国家航空航天局气候模拟中心的超级计算机，戈达德太空飞行中心，2010，公有领域（维基共享资源）；NCEP--NCAR 再分析第 1 版的地表气温与大气顶辐射，数据由美国科罗拉多州博尔德市 NOAA PSL 提供。
  \item 第~24~章（冰川学）：格陵兰西南部的一条出口冰川，NASA / Jim Yungel，冰桥行动（Operation IceBridge），2014，公有领域（维基共享资源）；格陵兰地表高程取自 ETOPO1 冰面数据，NOAA NCEI；北极海冰范围取自 NSIDC 海冰指数第 4 版（Fetterer 等，2025，doi:10.7265/a98x-0f50）。
  \item 第~25~章（遥感与地理信息系统）：俄勒冈州三姊妹火山中心 1995--2001 年的干涉图，美国地质调查局（C.~Wicks 等，2002），公有领域（维基共享资源）；实验室光谱取自 ECOSTRESS 光谱库第 1.0 版（S.~K.~Meerdink、S.~J.~Hook、D.~A.~Roberts 与 E.~A.~Abbott，2019；A.~M.~Baldridge 等，2009），经美国加利福尼亚州帕萨迪纳市加州理工学院喷气推进实验室惠允转载。
  \item 第~26~章（行星内部与表面）：亚利桑那州的巴林格陨石坑，USGS/D.~Roddy，公有领域（维基共享资源）；奥林帕斯山，维京号轨道器镶嵌图，NASA/JPL/USGS，公有领域（维基共享资源）。撞击坑产出与年代函数取自 G.~G.~Michael 的 Craterstats 软件，依据 G.~Neukum、B.~A.~Ivanov 与 W.~K.~Hartmann（2001）。
  \item 第~27~章（行星大气与宜居性）：火星上一个干涸河谷网络，火星快车 HRSC，ESA/DLR/FU Berlin，CC BY-SA 3.0 IGO（维基共享资源）；土卫六上的利吉亚海，卡西尼雷达镶嵌图，NASA/JPL/USGS，公有领域（维基共享资源）。宜居带拟合取自 R.~K.~Kopparapu 等（2013）；甲烷蒸气压取自 NIST 化学网络手册。
  \item 第~28~章（系外行星）：詹姆斯·韦伯太空望远镜拍摄的 HIP~65426~b，NASA/ESA/CSA，A~Carter（UCSC）、ERS~1386 团队与 A.~Pagan（STScI），CC BY~4.0（维基共享资源）。已确认行星取自 NASA 系外行星档案的行星系统综合参数表，访问日期 2026 年 10 月 23 日：本研究使用了 NASA 系外行星档案，该档案由美国加州理工学院运营，依据美国国家航空航天局系外行星探索计划的合同。质量--半径模型曲线取自 L.~Zeng 等（2016，2019）。
  \item 第~29~章（银河系与星系）：GigaGalaxy Zoom 项目的银河系全景，ESO/S.~Brunier，CC BY~4.0（维基共享资源）；旋涡星系 M101，欧洲航天局与 NASA（哈勃，heic0602；哈勃数据 K.~D.~Kuntz、F.~Bresolin、J.~Trauger、J.~Mould 与 Y.-H.~Chu；图像处理 Davide De~Martin），CC BY~4.0（维基共享资源）。银河系的圆周速度取自 A.-C.~Eilers、D.~W.~Hogg、H.-W.~Rix 与 M.~K.~Ness（2019）；质量模型取自 E.~Pouliasis、P.~Di~Matteo 与 M.~Haywood（2017）；大麦哲伦云的经典造父变星取自 OGLE，I.~Soszy\'nski 等（2015），Acta Astronomica 65, 297，由斯特拉斯堡 CDS 发布（VizieR J/AcA/65/297）。
  \item 第~30~章（元素从何而来）：蟹状星云，NASA、ESA、J.~Hester 与 A.~Loll（亚利桑那州立大学），公有领域（维基共享资源）；NGC~4993 中 GW170817 的千新星，NASA 与 ESA，致谢 A.~J.~Levan、N.~R.~Tanvir、A.~Fruchter 与 O.~Fox（哈勃，heic1717a），CC BY~4.0（维基共享资源）。太阳丰度取自 M.~Asplund、A.~M.~Amarsi 与 N.~Grevesse（2021）；中子俘获截面取自 KADoNiS v0.3（I.~Dillmann、R.~Plag、F.~K\"appeler 与 T.~Rauscher）；半衰期取自 IAEA 核素图表（ENSDF）；Geneva--Copenhagen 巡天的金属丰度取自 L.~Casagrande 等（2011）的再分析，由斯特拉斯堡 CDS 发布（VizieR J/A+A/530/A138）。
\end{itemize}
